\documentclass[../Main.tex]{subfiles}


\begin{document}

\section{Einleitung}

\subsection{Motivation und Relevanz}

Diese Maturitätsarbeit entstand aus dem Wunsch, nicht nur Spiele zu spielen, sondern sie genauer zu verstehen und selbst ein solches Spielerlebnis zu erschaffen.

Das Hack-and-Slash-Genre stellt aus unserer Sicht besonders hohe Anforderungen an das Zusammenspiel einzelner Spielsysteme, da Spielerbewegung, Kampfsystem, Gegnerverhalten und Leveldesign in Echtzeit ineinandergreifen. Im Zentrum unseres Interesses stehen dabei insbesondere das Gegnerverhalten und die Gegnerpositionierung.

In unserer Wahrnehmung erschienen Genrevertreter wie \game{Ninja Gaiden Sigma Plus} oder \game{God of War} in einer Phase mit vergleichsweise häufigen Neuveröffentlichungen des Genres, während heute nur noch wenige grosse Titel erscheinen. Dadurch existiert einerseits ausreichend Referenzmaterial, um zentrale Genremerkmale herauszuarbeiten, andererseits sind wir bei der Umsetzung nicht an aktuelle Genrekonventionen gebunden, wodurch sich zusätzlicher Freiraum für eigene, moderne Interpretationen klassischer Mechaniken ergibt.

\subsection{Erkenntnisfrage} \label{sec:Erkenntnisfrage}

Im Zentrum unserer Arbeit steht folgende Erkenntnisfrage:

\vspace{\baselineskip}

\begin{quote}
\textit{Wie müssen Gegnerverhalten und Gegnerplatzierung in einem 3D-Hack-and-Slash-Prototyp konzipiert sein, damit ein zusammenhängender Kampfabschnitt mit nachvollziehbarem Spielfluss, angemessener Schwierigkeit und befriedigendem Spielgefühl entsteht?}
\end{quote}

\vspace{\baselineskip}

Diese Frage setzt sich aus zwei inhaltlichen Komponenten zusammen. Zum einen geht es um das Gegnerverhalten und das Kampfgefühl, das in Kapitel~\ref{sec:Referenzspielanalyse} anhand einer Analyse bestehender Referenzspiele untersucht wird. Zum anderen geht es um Gegnerplatzierung und Leveldesign, dessen Prinzipien in Kapitel~\ref{sec:Leveldesign} literaturbasiert erarbeitet und anschliessend auf die Gestaltung unseres Kampfabschnitts angewendet werden. Beide Komponenten werden dabei in Kapitel~\ref{sec:Evaluation} in Bezug auf Spielfluss, Schwierigkeit und Spielgefühl evaluiert.

\subsection{Zielsetzung und Produktdefinition}

\textcolor{red}{Am Ende den Vertrag, das Spiel und den Text unten Abgleichen}

Ziel dieser Arbeit ist die Entwicklung eines spielbaren 3D-Hack-and-Slash-Prototyps mit der Game Engine Godot. Wodurch unsere Erkenntnisfrage geklärt werden soll.

\textcolor{red}{stuff den wir gemacht haben.}

\begin{itemize}[leftmargin=1.5cm]
    \item \textcolor{red}{Abgrenzung: Was steht in diesem Begleitbericht, was in der technischen Dokumentation? Titel des anderen Dokuments nennen und darauf verweisen.}
    \item \textcolor{red}{Produktbestandteile laut Vertrag (A.4) aufzählen: Prototyp, Kampfabschnitt, mind. zwei Gegnertypen, Architektur-Doku, Analyse, Playtests.}
\end{itemize}



\subsection{Aufbau der Arbeit} \label{sec:Aufbau}

\begin{itemize}[leftmargin=1.5cm]
    \item \textcolor{red}{Den Bogen der Arbeit ausformulieren: Frage $\rightarrow$ Grundlagen und Methodik (Kap.~\ref{sec:Grundlagen}) $\rightarrow$ von den Vorbildern lernen (Kap.~\ref{sec:Referenzspielanalyse}) $\rightarrow$ aus der Literatur lernen (Kap.~\ref{sec:Leveldesign}) $\rightarrow$ Vorschläge ableiten und bauen (Kap.~\ref{sec:Umsetzung}) $\rightarrow$ testen (Kap.~\ref{sec:Evaluation}) $\rightarrow$ Antwort (Kap.~\ref{sec:Schlussfolgerungen}) $\rightarrow$ Rückblick (Kap.~\ref{sec:Reflexion}).}
    \item \textcolor{red}{Pro Kapitel 1--2 Sätze: was passiert dort und \emph{warum steht es an dieser Stelle} (Merkblatt: Aufbau erklären und begründen).}
    \item \textcolor{red}{Ein Satz zur Zweiteilung: Begleitbericht und technische Dokumentation.}
    \item \textcolor{red}{Umfang ca.\ eine halbe Seite. Erst ganz am Schluss schreiben, wenn alle Kapitel stehen.}
\end{itemize}

\newpage

\end{document}